\subsection*{Justificacion Fila 9:}

\[\pi_D(S) \bowtie R\]

Primero se proyecta la columna $D$ de $S$. Eliminamos los renglones duplicados, eliminamos un $3$ ya que aparece dos veces.

\begin{center}
$\pi_{D}(S) =$ \quad
\begin{tabular}{|c|}
\hline
\textbf{D} \\ \hline
  3 \\
  1 \\
  0 \\
  2 \\ \hline
\end{tabular}
\end{center}

\begin{center}
\begin{minipage}{0.35\textwidth}
  \centering
  $R =$ \quad
  \begin{tabular}{|c|c|}
    \hline
    \textbf{A} & \textbf{B} \\ \hline
    1 & x \\
    2 & y \\
    2 & z \\
    3 & x \\
    9 & a \\ \hline
  \end{tabular}
\end{minipage}
\end{center}

$\pi_{D}(S) $ no comparte atributos con $R$, así que el join natural equivale al producto cartesiano.
\[4 \times 5 = \text{\textbf{20}}\]